% TOP_PROBLEM: {"id": 3000038, "title": "Generic rigidity in three dimensions", "category_id": 2, "category": {"id": 2, "name": "combinatorics", "display_name": "Combinatorics", "description": "Counting problems, graph theory, discrete structures.", "slug": "combinatorics", "order_index": 2, "created_at": "2026-07-31T15:26:25.671Z"}, "status": "partially_solved", "problem_number": "AMR-029-0038", "set_id": 13, "statusQualification": "The deterministic target is distinguished from known randomized symbolic-rank tests."}
% FIRST_POSTED: 2026-10-01
% ENTRY_KIND: statement_only
% UnsolvedMath ID 3000038; source catalogue code AMR-029-0038
% !TeX program = lualatex
% Definition and sources only; this file is not an LLM solution attempt.
\documentclass[11pt,a4paper]{article}
\usepackage[margin=25mm]{geometry}
\usepackage{fontspec}
\setmainfont{lmroman10-regular.otf}[BoldFont=lmroman10-bold.otf,
 ItalicFont=lmroman10-italic.otf,BoldItalicFont=lmroman10-bolditalic.otf]
\usepackage{amsmath,amssymb,mathtools}
\usepackage{unicode-math}
\setmathfont{latinmodern-math.otf}
\AtBeginDocument{\let\setminus\smallsetminus}
\usepackage{xurl,hyperref}
\hypersetup{colorlinks=true,urlcolor=blue}
\setcounter{secnumdepth}{0}
\setlength{\parindent}{0pt}
\setlength{\parskip}{5pt}
\setlength{\emergencystretch}{3em}
\begin{document}
\section*{Generic rigidity in three dimensions}\label{3000038}
{\small UnsolvedMath ID 3000038 · AMR-029-0038 · Combinatorics · AMR Open Problem Lists}
\subsection{Definitions and mathematical statement}
For a finite simple graph $G=(V,E)$ and placement $p:V\to\mathbb R^3$,
the rigidity matrix $R_3(G,p)$ has three columns per vertex and one row
per edge $uv$. That row contains $p(u)-p(v)$ in the columns of $u$,
$p(v)-p(u)$ in those of $v$, and zero elsewhere. Coordinates are generic
when algebraically independent over $\mathbb Q$; the resulting rank is
then a graph invariant.

For $n=|V|\ge3$, generic rigidity in three dimensions is equivalent to
\[
                           \operatorname{rank}R_3(G,p)=3n-6.
\]
For $n\le2$, use the same geometric definition: every sufficiently
nearby equivalent realization must be congruent; these small cases
are directly decidable.

Can generic rigidity be recognized in deterministic polynomial time
from the graph alone? Seek a combinatorial rank characterization or
another deterministic method, beyond randomized evaluation of the
symbolic matrix, which the source already records as available.
Infinitesimal motions are vectors in the kernel of $R_3$; the allowed
trivial motions come from translations and rotations. The question
concerns generic placements, not accidental rank drops at special
coordinates. It also differs from global rigidity: a graph may resist
all local flexes while admitting a distant, noncongruent realization
with exactly the same edge lengths.
\subsection{Short English statement}
Can one deterministically recognize generic three-dimensional
rigidity in polynomial time from the graph?
\subsection{Sources}
\begin{itemize}
\item[S1] ULAM.AI and the UnsolvedMath Contributors, supplied problems.json, record 3000038 (AMR-029-0038), AMR Open Problem Lists. Catalogue identity and source transcription; CC BY 4.0.
\url{https://huggingface.co/datasets/ulamai/UnsolvedMath}.
\item[S2] Egres Open - List of open problems, Open-problem page: Generic rigidity in three dimensions. Original source indexed by AMR-029-0038.
\url{https://oldlemon.cs.elte.hu/egres/open/List_of_open_problems}.
\item[S3] EGRES Open, Generic rigidity in three dimensions. Individual source page and explanatory remarks.
\url{https://lemon.cs.elte.hu/egres/open/Generic_rigidity_in_three_dimensions}.
\end{itemize}
\end{document}
